%---------- Inleiding ---------------------------------------------------------

\section{Inleiding}%
\label{sec:inleiding}

België heeft een militair engagement met de NAVO waarbij geplande en of gevraagde middelen ingezet worden binnen een coalitie van verschillende naties tijdens militaire operaties. Om een militaire operatie te doen slagen is een gezamenlijk en efficiënt communicatiesysteem nodig. Dit systeem dient om de verschillende verzamelde inlichtingen te delen met de coalitieleden. De gegevens leiden tot verschillende militaire interventies die een gezamenlijk effect creëren op het slagveld.

Echter, dat communicatiesysteem opbouwen en onderhouden is voornamelijk de taak van één natie. De reden is eenvoudig, elke natie is onafhankelijk om te beslissen welke communicatiemiddelen hij aankoopt. Dat wil zeggen dat de NAVO de naties niet zal verplichten tot een bepaalde aankoop van vendor en of technologieën. Het is vanzelfsprekend dat dit een grote impact heeft op de snelheid waarbij gegevens gedeeld worden en dan nog niet gesproken over de trage opbouw van zo’n communicatiesysteem. 

De langdurige interventie in Afghanistan heeft veel van deze beperkingen blootgelegd:
\begin{itemize}
    \item Snel delen van inlichtingen tussen verschillende partners is belangrijk.
    \item Data is gevoelig en moet beschermd worden.
\end{itemize}

De combinatie van deze twee vraagt heel wat denkwerk. Bijgevolg heeft de NAVO het FMN project opgestart. FMN heeft als doel om:
\begin{itemize}
    \item Te standaardiseren onafhankelijk van vendor.
    \item Zero-day interoperability.
    \item Verschillende beveiligde gegevensdomeinen aanbieden en onderling verbinden.
\end{itemize}

\subsection{Probleemstelling}%
Het grote aantal verschillende netwerken leidt tot beperkingen in interoperabiliteit en vertraagt de inzetbaarheid van militaire eenheden in coalitiemissies. Federated Mission Networking (FMN) biedt een potentieel antwoord op deze problemen door een ‘tijdelijk’ missie-netwerk te ontwikkelen dat landen in staat stelt om tijdens missies snel en veilig gegevens te delen over dit missie-netwerk. Dit onderzoek richt zich op de integratie van België in het Federated Mission Networking project door gebruik te maken van de nieuwe nationale SD-WAN technologieën.
%---------- Stand van zaken ---------------------------------------------------
\subsection{Centrale onderzoeksvraag}%
Hoe kan België optimaal profiteren van deelname aan het FMN-initiatief om interoperabiliteit in multinationale missies te verbeteren, na de instap in de SD-WAN technologie.
\subsubsection{Deelvragen met doelstellingen}%
\paragraph{Welke technische en organisatorische vereisten zijn nodig voor België om succesvol deel te nemen aan FMN?:\\}
\textbf{Doelstelling}: Identificeren van de specifieke technische infrastructuren en organisatorische middelen die nodig zijn voor Belgische deelname aan FMN, waaronder netwerken, systemen en beveiligingsprotocollen.

\textbf{Doelstelling}: Welke netwerk rol(len) kan België vervullen binnen FMN?
\paragraph{Wat zijn de voordelen van het implementeren van een permanente FMN-infrastructuur in België:\\}
\textbf{Doelstelling}: Onderzoeken hoe een permanente FMN-infrastructuur de reactietijd en operationele effectiviteit van Belgische eenheden binnen coalitiemissies kan verbeteren, met focus op snelheid en interoperabiliteit.
\paragraph{Welke stappen zijn nodig om een werkend labomodel te creëren voor het testen en valideren van FMN?:\\}
\textbf{Doelstelling}: Een gebruiksconcept ontwikkelen voor een testomgeving die voldoet aan de vereisten van 'Spiral 4'.

\section{Literatuurstudie}%
\label{sec:literatuurstudie}
\subsection{Inleiding}
Federated Mission Networking (FMN) is een initiatief van de NAVO dat streeft naar verbeterde interoperabiliteit tussen bondgenoten en partners tijdens gezamenlijke militaire operaties. In de essentie gaat het over een vast kader creëren waar elke natie met zijn eigen hardware op kan aansluiten om vervolgens informatie met elkaar uit te wisselen.
 
\subsection{Achtergrond en ontwikkeling}
FMN is gebaseerd op lessen uit het Afghanistan Mission Network ~\autocite{FMN} en richt zich op het verbinden van NAVO-landen en partners met een netwerk dat kan worden aangepast aan operationele eisen. Het initiatief omvat een technische architectuur, best practices en gemeenschappelijke standaarden ~\autocite{FMN_Guide}. Deze conventies dienen door alle partijen toegepast te worden.

\subsection{Software Defined Networking (SDN) en beveiliging via IPsec}
In een traditionele netwerkomgeving heeft elk apparaat, zoals routers en switches, zijn eigen controlemechanismen. SDN scheidt echter de beheerslaag (control plane) van de gegevenslaag (data plane), waardoor een centrale controller het netwerk beheert en stuurt. Dit maakt het mogelijk om snel aanpassingen te doen op basis van de behoeften van applicaties en diensten ~\autocite{SDN}.

\subsubsection{Beveiliging en Interoperabiliteit: GRE over IPsec}
FMN maakt gebruik van GRE over IPSEC om de verschillende naties met elkaar te verbinden.
GRE kan verschillende soorten verkeer encapsuleren, zoals unicast-, multicast-, broadcast- en MPLS-verkeer. GRE biedt echter geen bescherming voor de overgedragen payload.
Internet Protocol Security (IPsec) daarentegen zorgt voor vertrouwelijkheid, integriteit en authenticatie van de payloads die via IPsec-tunnels worden verzonden. IPsec werkt echter uitsluitend met IP-pakketten.
De GRE-over-IPsec combineert de flexibiliteit van GRE met de beveiliging van IPsec. GRE encapsuleert de pakketten, terwijl IPsec deze pakketten versleutelt en transporteert via een IPsec-tunnel ~\autocite{GRE}.

\subsubsection{Uitdaging binnen SDN}
Belgische Defensie is vandaag ingestapt in de wereld van SDN (Software Defined Networking). De verschillende leveranciers die deze netwerktechnologie verkopen staan haaks op de idee FMN. Leveranciers hebben er belang bij om aan “vendor lock-in” te doen. Dat wil zeggen dat bedrijven die vandaag een keuze maken in technologie zich tevens vastbinden aan een bepaald ecosysteem. Je ben als het ware verplicht om alle functionele onderdelen van het SDN-systeem aan te kopen (controller, licenties etc.) van diezelfde leverancier. Dit staat haaks op de filosofie van FMN die juist gericht is op samenwerking tussen naties en het gebruik van gestandaardiseerde technologieën.

\subsubsection{Internet Key Exchange (IKE)}
IKE is een protocol dat gebruikt wordt om een beveiligde communicatieomgeving op te zetten door cryptografische sleutels uit te wisselen en  Security Associations (SA's) te configureren. Het is een essentieel onderdeel van het IPsec-protocolsuite, dat zorgt voor veilige communicatie over niet-vertrouwde netwerken ~\autocite{IKE}.

\subsubsection{Samenhang tussen SDN en IPsec}
Hoewel SDN primair gericht is op het dynamisch beheer van netwerkresources, speelt beveiliging een cruciale rol in de implementatie, vooral in militaire toepassingen zoals bij FMN het geval is. IPsec kan als een beveiligingslaag binnen een SDN-omgeving worden geïntegreerd om de vertrouwelijkheid en integriteit van data te waarborgen. De centrale controller van SDN kan ervoor zorgen dat IPsec-tunnels dynamisch worden opgezet en aangepast op basis van de huidige netwerkbehoeften, wat leidt tot een efficiëntere en veiligere netwerkarchitectuur.

\subsection{Waar we vandaag staan}
Binnen de context van FMN verwijst een "spiral" naar een fase van incrementele ontwikkeling en implementatie in het project. Deze ‘spirals’ volgen elkaar op in opeenvolgende stappen, waarbij elke ‘spiral’ nieuwe specificaties, technische verbeteringen en operationele functionaliteiten toevoegt aan het netwerk. Het idee achter deze benadering is om telkens kleine, maar cruciale verbeteringen door te voeren die zorgen voor grotere interoperabiliteit tussen NAVO-leden en partners. Dit onderzoek zal zich afspelen in ‘Spiral 4’~\autocite{Spiral4}.

% Voor literatuurverwijzingen zijn er twee belangrijke commando's:
% \autocite{KEY} => (Auteur, jaartal) Gebruik dit als de naam van de auteur
%   geen onderdeel is van de zin.
% \textcite{KEY} => Auteur (jaartal)  Gebruik dit als de auteursnaam wel een
%   functie heeft in de zin (bv. ``Uit onderzoek door Doll & Hill (1954) bleek
%   ...'')

%---------- Methodologie ------------------------------------------------------
\section{Methodologie}%
\label{sec:methodologie}
\subsection{Verschillende fases}
\paragraph{Literatuurstudie en documentanalyse:\\} 
Door het analyseren van bestaande literatuur en NATO-richtlijnen wordt een solide theoretische basis gelegd. Dit biedt inzicht in de vereisten van FMN Spiral 4 en de rol van SD-WAN in het ondersteunen van interoperabiliteit. Deze fase identificeert de technische uitdagingen en biedt criteria voor de evaluatie van SD-WAN apparatuur.

\paragraph{Vergelijkende analyse:\\}
Op basis van de literatuur worden Cisco- en Aruba-apparaten vergeleken aan de hand van criteria zoals kostprijs, beheer, integratie en compatibiliteit met andere ecosystemen binnen FMN. Hierbij wordt onderzocht of één enkel apparaat alle vereiste netwerkrollen kan uitvoeren, wat cruciaal is voor eenvoud en efficiëntie.

\paragraph{Praktische implementatie van FMN:\\}
Een prototype van FMN Spiral 4 wordt opgezet met de geselecteerde SD-WAN-technologie. Dit gebeurt in een gecontroleerde testomgeving waarin typische FMN-scenario's worden gesimuleerd. De resultaten bieden inzicht in de technische en operationele haalbaarheid van de integratie.

\paragraph{Evaluatie en aanbevelingen:\\}
De resultaten van de implementatie en analyse worden geëvalueerd aan de hand van de centrale onderzoeksvraag. De aanbevelingen richten zich op de optimale rol van België binnen het FMN-initiatief en de implementatie van de benodigde technische infrastructuur.

\subsection{Tools}
\subsubsection*{Evaluatie van bruikbaarheid}
Om een weloverwogen keuze te kunnen maken in het type SD-WAN apparatuur toetsen we deze op basis van onderstaande criteria:
\begin{itemize}
    \item Kostprijs (initiële kosten, licentiekosten en onderhoudskosten)
    \item Beheer en gebruiksvriendelijkheid
    \item Integratie met bestaande militaire infrastructuur
    \item Capaciteit om alle netwerkrollen binnen FMN uit te voeren
\end{itemize}

\subsection{Technologiën}
\begin{itemize}
    \item Software-Defined Networking
    \item GRE-tunnel
    \item IPSEC
    \item Internet Key Exchange (IKE)
\end{itemize}

\subsection{Tijdsplanning en deliverables}
\paragraph{Week 1-3:} 
Uitgebreide literatuurstudie, kennis verbeteren over de verschillende technologieën gebruikt bij FMN.
Deliverables: Inleiding, uitgebreid literatuuroverzicht, methodologie, overzicht van vereisten en criteria om aan FMN te kunnen deelnemen.
\paragraph{Week 4-8:}
Praktische implementatie van FMN en testen.
Deliverables: gebruiksconcept en testresultaten.
\paragraph{Week 9-10:}
Validatie en evaluatie van de praktische implementatie.
Deliverables: evaluatierapport a.d.h.v. de deelvragen en doelstellingen van dit onderzoek 


%---------- Verwachte resultaten ----------------------------------------------
\section{Verwacht resultaat, conclusie}%
\label{sec:verwachte_resultaten}
De resultaten van dit onderzoek kunnen inzicht bieden in hoe een permanente inzet van FMN-infrastructuur in België mogelijk de reactietijd bij missies versnelt en de interoperabiliteit versterkt. Daarnaast wordt er een advies geformuleerd over het gebruik van SD-WAN apparatuur binnen de FMN-context volgens de NATO-richtlijnen. Hoewel de bevindingen wijzen op het potentieel van FMN om, met een gedeelde Cloud-omgeving, coalitiemissies efficiënter en veiliger te maken, benadrukken ze ook de noodzaak van verder onderzoek en evaluatie. Deze inzichten zijn van cruciaal belang voor het werkveld, waar effectieve communicatie tussen coalitiepartners steeds belangrijker wordt.